% SPDX-License-Identifier: MIT
% Copyright (c) 2026 handsomeZR-netizen
%
% Original CUMCM 2026 electronic-paper assembly template for mathmodel-skill.
% This is a rules-aligned repository template, not an official contest template.
% It intentionally omits the commitment form, numbering page, and identity data.

\documentclass[UTF8,12pt,a4paper,fontset=fandol]{ctexart}

\usepackage[margin=2.5cm]{geometry}
\usepackage{amsmath,amssymb,amsthm}
\usepackage{graphicx}
\usepackage{float}

% 默认图片就地排版 [H]：图片紧跟正文引用位置，避免浮动漂离产生空页/孤图
\makeatletter
\def\fps@figure{H}
\makeatother
\usepackage{booktabs}
\usepackage{longtable}
\usepackage{multirow}
\usepackage{listings}
\usepackage{xcolor}
\usepackage{hyperref}
\usepackage{url}
\usepackage{cite}
\usepackage{caption}
\usepackage{subcaption}
\usepackage{enumitem}

% pandoc 生成 longtable 表格所需的 calc/array 与 none 计数器
\usepackage{array}
\usepackage{calc}
\newcounter{none}

\hypersetup{
  unicode=true,
  hidelinks,
  pdfauthor={},
  pdftitle={},
  pdfsubject={},
  pdfkeywords={}
}
\urlstyle{same}

% ---- 定理/引理/推论/定义环境 (amsthm, 证明型题目用) ----
\newtheorem{theorem}{定理}[section]
\newtheorem{lemma}[theorem]{引理}
\newtheorem{corollary}[theorem]{推论}
\newtheorem{definition}[theorem]{定义}

\setlength{\parindent}{2em}
\setlength{\parskip}{0pt}
\linespread{1.25}
\captionsetup{font=small,labelfont=bf}

% 允许版心内行界适度伸缩，消除长公式/长词造成的行内数学 overfull(右边界溢出)。
% 配合行内公式内的 \allowbreak，保证所有文字与公式左右边界与版心对齐。
\setlength{\emergencystretch}{3em}

% 章节标题层级：一级"一、"黑体居中(字号大于正文)，二级"1.1"黑体左对齐
\ctexset{
  section = {
    format = {\centering\zihao{3}\heiti},
    name = {},
    number = \chinese{section},
    aftername = {、},
  },
  subsection = {
    format = {\raggedright\zihao{4}\heiti},
    name = {},
  },
  subsubsection = {
    format = {\raggedright\zihao{-4}\heiti},
    name = {},
  },
}
\lstset{
  basicstyle=\small\ttfamily,
  breaklines=true,
  frame=single,
  numbers=left,
  numberstyle=\tiny,
  showstringspaces=false,
  columns=fullflexible,
  keepspaces=true
}

% Pandoc may emit \tightlist in a LaTeX fragment.
\providecommand{\tightlist}{%
  \setlength{\itemsep}{0pt}\setlength{\parskip}{0pt}}

% These counters enforce the electronic-paper structure during compilation.
% The abstract occupies page 1. The body begins on page 2 and is limited to
% 30 pages; references and the optional no-AI declaration count as body pages,
% while the appendix does not.
\newcounter{mathmodelbodystartpage}
\newcounter{mathmodelbodypages}

\begin{document}

% MATHMODEL_* fields are fail-closed rendering tokens. render_paper.py fills
% them from decision_log.paper_metadata or explicit CLI arguments.
\pagenumbering{arabic}
\pagestyle{plain}
\thispagestyle{plain}

\begin{center}
  {\zihao{3}\bfseries MATHMODEL_TITLE}\par
  \vspace{1em}
  {\zihao{4}\bfseries 摘要}\par
\end{center}

% center(trivlist) 结束后其内部 \par 会吞掉下一段的缩进 token，
% 导致摘要正文第一段无首行缩进。此处用 \indent 强制该段缩进 2 字。
\indent
% MATHMODEL:SECTION abstract

\par\noindent\textbf{关键词：}MATHMODEL_KEYWORDS

% The electronic paper must start with a one-page abstract sheet. Starting the
% body on any page other than page 2 is a submission-format error.
\clearpage
\ifnum\value{page}=2\relax\else
  \PackageError{mathmodel-mathmodel}
    {CUMCM abstract must fit on the first page}
    {Shorten the title, abstract, or keywords before submission.}
\fi
\setcounter{mathmodelbodystartpage}{\value{page}}

% The Markdown section files own their headings. No table of contents is added.
% MATHMODEL:SECTION 1_problem_restate

% MATHMODEL:SECTION 2_problem_analysis

% MATHMODEL:SECTION 3_assumptions

% MATHMODEL:SECTION 4_notation

% MATHMODEL:SECTION 5_models

% MATHMODEL:SECTION 6_sensitivity

% MATHMODEL:SECTION 7_evaluation

% MATHMODEL:SECTION 8_references

% An explicit no-AI declaration belongs immediately after the references.
% When AI was used, render_ai_usage.py does not create this source file and the
% renderer removes this whole optional block.
% MATHMODEL:OPTIONAL mathmodel_no_ai_statement BEGIN
\par\bigskip
% MATHMODEL:SECTION mathmodel_no_ai_statement
% MATHMODEL:OPTIONAL mathmodel_no_ai_statement END

% Flush deferred body floats before measuring. At this point \value{page} is
% the first appendix page, so subtracting the body start yields the body length.
\clearpage
\setcounter{mathmodelbodypages}{%
  \numexpr\value{page}-\value{mathmodelbodystartpage}\relax}
\typeout{MATHMODEL:CUMCM_BODY_PAGES=\arabic{mathmodelbodypages}}
\ifnum\value{mathmodelbodypages}>30\relax
  \PackageError{mathmodel-mathmodel}
    {CUMCM body exceeds 30 pages}
    {Shorten the main text. Appendices are counted separately.}
\fi

\appendix
% MATHMODEL:SECTION appendix_code

\end{document}
